% sec_negative.tex -- negative-results sections. \input into a document using preamble.tex.
% Label prefix: neg:. Sources: worker reports and chains listed in report.md of the writer.

\section{Routes to Hypothesis (A) that fail}\label{sec:routes}

Five routes to Hypothesis~(A) were written out as chains of steps. Each step was tagged
proved, cited, claimed or missing, and the first unproved step was tested. This section
states, for each route, what is proved and where the route fails. All programs were written
by workers of one AI model family; no program was run by anyone outside it.

We use the following status words for evidence. \emph{Computed exactly}: exact arithmetic or
a closed form. \emph{Computed in floating point}: a deterministic computation in double
precision. \emph{Monte Carlo by two separate programs}: two codes written independently,
neither copied from the other. \emph{Monte Carlo by one program}: one code, possibly rerun.
\emph{Cited, as stated in its abstract}: the source was read at abstract or search-snippet
level only. Numerical evidence is never proof. Where the sources do not give seeds, we say so.

Two discriminators are used throughout. Compact $\U(1)$ in four dimensions has no gap at weak
coupling \cite{Guth80,FS82} (both as stated in their abstracts). Finite gauge groups in four
dimensions have exponential decay of correlations at weak coupling \cite{AC22}. A route to (A)
that applies verbatim to $\U(1)$ proves something false. A route that cannot reach the finite
groups is incomplete, not wrong.

Throughout, $\beta$ is the Wilson coupling: the weight is
$\exp\bigl(\tfrac{\beta}{2}\sum_p \operatorname{Re}\Tr U_p\bigr)$ against Haar measure on
$\SU(2)^E$. Boxes are $\Lambda$; boundary conditions are $\tau$.

\subsection{Positive curvature}\label{neg:sub-curv}

\emph{Mechanism.} The Wilson measure is $e^{-V}\,d\mathrm{Haar}$ on the Riemannian manifold
$\SU(2)^E$, with $V=-\frac{\beta}{2}\sum_p\operatorname{Re}\Tr Q_p$. If the Bakry--\'Emery tensor
$\mathrm{Ric}+\operatorname{Hess}V$ is bounded below by $K>0$, a Poincar\'e inequality holds with
constant $1/K$, uniformly in $\Lambda$ and $\tau$. Uniform Poincar\'e inequalities give decay of
correlations at a rate independent of the supports, and such decay gives a transfer-matrix gap.

Fix the metric $\langle X,Y\rangle=-\Tr(XY)$ on $\mathfrak{su}(2)$. Then $\mathrm{Ric}=1$
\cite{SZZ23}. Let $B(U)=1+\operatorname{Hess}V(U)$ on $T_U\SU(2)^E$. Let $H_U$ be the orthogonal
complement of the tangent space of the gauge orbit through $U$. Put
$\kappa_H(U)=\min\{B(U)(v,v):v\in H_U,\ |v|=1\}$.

\begin{lemma}\label{neg:lem-horizontal}
Let $f$ be a smooth gauge-invariant function on $\SU(2)^{E(\Lambda)}$. Then $\nabla f(U)\in H_U$
for every $U$, and the Langevin semigroup $P_t$ of $\mu_{\Lambda,\tau}$ maps gauge-invariant
functions to gauge-invariant functions. Consequently, if $\kappa_H(U)\ge K>0$ for every $U$,
then $\Var_{\mu_{\Lambda,\tau}}(f)\le K^{-1}\int|\nabla f|^2\,d\mu_{\Lambda,\tau}$ for every
gauge-invariant $f$, with $K$ independent of $\Lambda$ and $\tau$.
\end{lemma}

\begin{proof}
The gauge group acts by left and right translations in each factor. The metric is
bi-invariant, and $V$ and Haar measure are invariant. So the generator commutes with the
action, and $P_t$ preserves invariant functions. The differential of an invariant $f$
vanishes on orbit directions, so $\nabla f\in H_U$. Bochner's formula gives
$\Gamma_2(f)=\|\nabla^2 f\|^2+B(\nabla f,\nabla f)\ge\kappa_H\Gamma(f)$ along $P_tf$, and the
Bakry--\'Emery argument \cite[Thm.~1.4 and (4.7)]{SZZ23} gives the inequality.
\end{proof}

The gain over the full tensor is not small. On the periodic $3^4$ lattice the minimum of $B$
over all directions was $-6.4$, $-9.0$, $-12.2$ at $\beta=12,24,48$, against $\kappa_H=-1.6$,
$-2.0$, $-1.5$ (Monte Carlo by one program). The lemma was rechecked by a second worker by
argument only.

\begin{proposition}\label{neg:prop-flat}
\textup{(i)} At every flat configuration ($Q_p=1$ for all $p$), $\kappa_H=1$ on the periodic torus.
\textup{(ii)} For one plaquette with $Q=-1$, $B$ has eigenvalues $1-2\beta$ with multiplicity $3$,
on horizontal vectors, and $1$ with multiplicity $9$. Hence a pointwise bound
$\kappa_H\ge K>0$ fails for the one-plaquette system when $\beta\ge 1/2$.
\textup{(iii)} For compact $\U(1)$ the constant in \textup{(i)} is $0$.
\end{proposition}

\begin{proof}
(i) At $Q_p=1$, $\operatorname{Hess}V(t,t)=\frac{\beta}{2}\sum_p|Y_p(t)|^2\ge0$, where $Y_p(t)$
is the signed sum of the transported generators along $p$. So $B\ge1$, and $B=1$ exactly on
horizontal zero modes of the curl form. On the torus these exist (the twelve constant modes).
(ii) At $Q=-1$ the sign of the second-order term reverses:
$\operatorname{Hess}V(t,t)=-\frac{\beta}{2}|Y(t)|^2$. On $\mathfrak{su}(2)^4$ the form $|Y|^2$ has
eigenvalue $4$ with multiplicity $3$ and $0$ with multiplicity $9$. The gauge orbit through a
configuration with $Q=-1$ central has dimension $9$ and lies in the kernel. So the three
eigenvalues $1-2\beta$ are horizontal.
(iii) $\U(1)$ is flat, so $\mathrm{Ric}=0$.
\end{proof}

Part (ii) was computed exactly and confirmed by finite differences in three separate programs
(builder, challenger, and the builder's child). Part (i) is the one place in the route where
non-abelian structure enters. The cited criterion $\kappa_H\ge1-12\beta$ \cite[(4.2), (4.8),
Thm.~1.4]{SZZ23} gives (A) for $\beta<1/12$. On the four-dimensional periodic $3^4$ lattice,
the configuration with all links $1$ except one link $-1$ gives $\kappa_H=1-4.31\beta$ for
$\beta=1,\dots,6$ (computed in floating point by one program).

The route needs a step that carries the curvature to large $\beta$. The step written was:

\begin{hypothesis}[Step S6]\label{neg:hyp-S6}
For large $\beta$, $\kappa_H(U)\ge 1/2$ except on configurations containing plaquettes with
$\operatorname{Re}\Tr Q_p<0$; such configurations have density $e^{-c\beta}$, so a local
perturbation argument applies.
\end{hypothesis}

\begin{observation}\label{neg:obs-curv}
Step S6 is false on every periodic box tested. On the periodic lattices $3^4$ and $4^4$ at
$\beta\in\{4,6,12,24,48\}$, $\kappa_H$ of equilibrium configurations lies in
$[-2.30,-1.12]$ (Table~\ref{neg:tab-curv}). No sample of either program has $\kappa_H\ge1/2$:
$0$ of $32$ (builder) and $0$ of $24$ (challenger). At $\beta=4$ and $6$ the fraction of
plaquettes with $\operatorname{Re}\Tr Q_p<0$ is $0.002$ and $0.000$. Status: Monte Carlo by two
separate programs. The sources do not give the seeds. Thermalization was short ($30+4k$
sweeps in the builder's program; cold start plus $300$ sweeps in the challenger's).
\end{observation}

\begin{table}[ht]
\centering
\small
\begin{tabular}{@{}lllll@{}}
\toprule
program & box & $\beta$ & samples & $\kappa_H$ \\
\midrule
builder & $3^4$ periodic & $4$, $6$ & $8$ each & $-2.2$ to $-1.5$ \\
builder & $3^4$ periodic & $12$, $24$, $48$ & $4$ each & $-2.30$ to $-1.17$ \\
builder & $4^4$ periodic & $4$ & $2$ & $-1.60$, $-1.88$ \\
builder & $4^4$ periodic & $6$ & $2$ & $-1.53$, $-1.58$ \\
challenger & $3^4$ periodic & $4$ & & $-1.79$ to $-1.46$ \\
challenger & $3^4$ periodic & $6$ & & $-1.96$ to $-1.39$ \\
challenger & $3^4$ periodic & $12$ & & $-1.52$ to $-1.12$ \\
challenger & $3^4$ periodic & $24$ & & $-1.93$ to $-1.12$ \\
challenger & $3^4$ periodic & $48$ & & $-2.04$ to $-1.55$ \\
challenger & $4^4$ periodic & $6$ & $2$ & $-1.77$, $-1.26$ \\
challenger & $4^4$ periodic & $24$ & $2$ & $-1.15$, $-1.67$ \\
\midrule
challenger & $4^4$ open & $4$ & & $-2.40$ to $-1.67$ \\
challenger & $4^4$ open & $6$ & $3$ & $+0.16$, $-0.32$, $-0.04$ \\
challenger & $4^4$ open & $12$, $24$, $48$ & & $+3$, $+9.6$, $+23$ \\
\bottomrule
\end{tabular}
\caption{Horizontal curvature $\kappa_H$ of equilibrium configurations, $\SU(2)$. The
challenger ran $20$ periodic samples on $3^4$ in total. The builder's own summary row gives
$[-2.3,-1.2]$ and the director's ledger gives $-1.1$ to $-2.3$; the union of the per-run values
of both programs is $[-2.30,-1.12]$.}
\label{neg:tab-curv}
\end{table}

The challenger found the mechanism, measured only. On the torus the negative directions are
the constant (toron) modes: the lowest eigenvector has zero-momentum weight $0.96$ to $0.99$ in
Landau gauge. On a one-loop ensemble without toron fluctuations, $\kappa_H\to\mathrm{Ric}-k\,c$
with $c=\beta\varepsilon^2$ and $c=1$ at equilibrium; $k=2.62$ and $2.51$ for $\SU(2)$,
constant for $\beta$ from $16$ to $1024$, so the threshold $c_0=\mathrm{Ric}/k$ is $0.38$ to
$0.40$. For $\SU(3)$, $k=3.42$ and $c_0=0.44$ (one seed). In the open $4^4$ box at $\beta\ge12$
the typical curvature is positive (Table~\ref{neg:tab-curv}). This does not rescue the route.
The criterion is pointwise, the open-box positivity grows like $\beta p_{\min}^2$, and every box
of side $L\gtrsim\sqrt\beta$ or with a periodic direction has soft modes; the last statement is
an extrapolation, not a measurement.

\begin{observation}\label{neg:obs-curv-missing}
The route needs, in place of S6, the following statement, which is missing. There are
$\ell,C<\infty$ and $m>0$ such that for every cube $B$ of side at least $\ell$, every pair of
boundary conditions $\tau,\tau'$ and every gauge-invariant $f$ supported in $B$ with
$\|f\|_\infty\le1$,
\[
|\mu_{B,\tau}(f)-\mu_{B,\tau'}(f)|\le C\sum_{e:\,\tau_e\ne\tau'_e}e^{-m\,\mathrm{dist}(\operatorname{supp}f,e)}.
\]
This is strong mixing, essentially the strong mass gap of \cite[Def.~2.3]{Ch21}. It is (A)
in a stronger form, so the route is circular. The two steps after it (strong mixing gives a
uniform log-Sobolev inequality; a uniform Poincar\'e inequality gives decay with a rate
independent of the supports) are claimed from memory, not read.
\end{observation}

\emph{Discriminators.} For $\U(1)$ the route fails at its first step, since $\mathrm{Ric}=0$;
for finite groups there is no Riemannian structure and the route does not apply.

\subsection{Center vortices}\label{neg:sub-vort}

\emph{Mechanism.} A twist of the coupling on a coclosed stack $\mathcal V$ of $L^2$ plaquettes
in the periodic box $(\Z/L\Z)^4$ defines the vortex free energy $F_L(\beta)=-\ln(Z^-/Z)$.
Reflection positivity turns a bound on $F_L$ into a bound on Wilson loops, so decay of $F_L$
like $e^{-\sigma L^2}$ gives an area law. The route then needs an area law to give a gap.

\begin{lemma}\label{neg:lem-Zminus}
For every $\beta>0$ and even $L$, $0<Z^-\le Z$; that is, $F_L\ge0$.
\end{lemma}

\begin{proof}
$Z^->0$ because its integrand is positive. Let $O[\mathcal V]$ be the ratio of twisted to
untwisted weights and $P=\frac12(1-O[\mathcal V])$. By the change of variables of
\cite[(4), (7)]{KT02}, $\langle P\rangle=\langle P\cdot\theta P\rangle$, where $\theta$ is the
reflection in a plane between sites that does not meet $\mathcal V$. Reflection positivity
\cite{OS78} (as stated in its abstract) gives $\langle P\cdot\theta P\rangle\ge0$. So
$\langle O[\mathcal V]\rangle\le1$.
\end{proof}

\begin{proposition}\label{neg:prop-arealaw}
Let $C$ be an $R\times T$ rectangle in the $(1,2)$ plane, with $L/R$ and $L/T$ powers of $2$
and the reflection planes placed as in \cite{KT02} (through sites, not bisecting $\mathcal V$).
Then $\langle W[C]\rangle_{\Lambda_L}\le2\bigl(\tfrac12(1-e^{-F_L})\bigr)^{RT/L^2}$
\cite[(11)]{KT02}. Consequently, if $F_L(\beta)\le KL^2e^{-\sigma L^2}$ for all $L=2^kL_0$, then
every limit point of the periodic states along $L=2^kL_0$ satisfies
$\langle W[C]\rangle\le2e^{-\sigma RT}$.
\end{proposition}

\begin{proof}
The inequality is \cite[(11)]{KT02}, read in full text by the worker; \cite{KT02} attribute it
to \cite{TY85}, which was not read. Since $F_L\ge0$, $\frac12(1-e^{-F_L})\le F_L/2$, so
$\ln\langle W\rangle_{\Lambda_L}\le\ln2+\frac{RT}{L^2}\bigl(\ln(KL^2/2)-\sigma L^2\bigr)\to\ln2-\sigma RT$.
$W$ is local, so the bound passes to weak limits.
\end{proof}

\begin{proposition}\label{neg:prop-noarea}
An area law does not imply a gap. Let $\mu$ be any lattice gauge measure with an area law and
let $\nu$ be the free massless lattice scalar field. Under $\mu\otimes\nu$ every Wilson loop has
the same expectation as under $\mu$, and the transfer matrix has no gap.
\end{proposition}

\begin{proof}
Wilson loops depend only on the gauge factor. The transfer matrix of the product is
$T_\mu\otimes T_\nu$, and the spectrum of its generator contains $0+\operatorname{spec}H_\nu$,
which accumulates at $0$ for the massless field.
\end{proof}

\begin{observation}\label{neg:obs-vort}
The route fails at two independent steps.
\textup{(1)} The decay of $F_L$ at every $\beta$, $\liminf_L(-\ln F_L)/L^2\ge\sigma(\beta)>0$, is
missing; it is the confinement problem. The step meant to approach it, $F_{L'}\le F_L$ for
$L\mid L'$, is claimed by a renormalization-group comparison that is exactly the step
\cite{IS07,IS08} object to in \cite{To07}, and is not closed. It is false for $\Z_2$ at
$\beta=0.6$ (Table~\ref{neg:tab-vort}). For $\SU(2)$ the data neither support nor refute it:
the errors are as large as the differences, and at $\beta=2.4$,
$F_4-F_3=+1.26\pm1.05$, the wrong sign at $1.2$ standard errors.
\textup{(2)} The implication ``area law $\Rightarrow$ gap in the zero-flux sector'' is missing
and false as an abstract implication (Proposition~\ref{neg:prop-noarea}); the theorem of
\cite{Ch21} goes the other way. The flux free energy controls the sectors of nonzero electric
flux, not the neutral sector.
Status: $\SU(2)$ values, Monte Carlo by one program (Metropolis with a step size adapted during
the run, so not strictly detailed balance); $\Z_2$ values, Monte Carlo by one program, rerun bit
for bit by a second worker; the three-dimensional $\Z_2$ check on $L=2$ was computed exactly by
two separate programs ($F(0.3)=0.0325242683$, $F(0.6)=0.6912790067$, agreement to $10^{-10}$).
The sources do not give the seeds.
\end{observation}

\begin{table}[ht]
\centering
\small
\begin{tabular}{@{}llllll@{}}
\toprule
group & $\beta$ & $F_2$ & $F_3$ & $F_4$ & status \\
\midrule
$\SU(2)$ & $1.8$ & $0.69\pm0.24$ & $-0.26\pm0.28$ & $0.49\pm0.40$ & MC, one program \\
$\SU(2)$ & $2.4$ & $3.06\pm0.42$ & $1.55\pm0.45$ & $2.81\pm0.95$ & MC, one program \\
$\Z_2$ & $0.3$ & $0.0774$ & $0.0041$ & $0.0061$ & MC, one program, rerun \\
$\Z_2$ & $0.6$ & $3.205$ & $8.244$ & $15.84$ & MC, one program, rerun \\
\bottomrule
\end{tabular}
\caption{Vortex free energy $F_L$ on $L^4$, $d=4$. $\Z_2$ errors: $0.002$ to $0.005$ at
$\beta=0.3$, $\pm0.25$ at $\beta=0.6$.}
\label{neg:tab-vort}
\end{table}

No proved or cited step of this chain uses non-abelian structure; Lemma~\ref{neg:lem-Zminus}
and Proposition~\ref{neg:prop-arealaw} hold verbatim for $\U(1)$ and $\Z_2$. The non-abelian
input would enter only through asymptotic freedom in the monotonicity step, which is claimed.
An area law at all couplings for a non-abelian group is known only in two dimensions
\cite[p.~8]{Ch18}.

\emph{Discriminators.} For $\U(1)$ the first failing step is the decay of $F_L$ (perimeter law
\cite{Guth80,FS82}, as stated in their abstracts); for $\Z_2$ at weak coupling the monotonicity
step fails while (A) is true there, so the route cannot reach the finite groups.

\subsection{A finite-size criterion}\label{neg:sub-fs}

\emph{Mechanism.} A Dobrushin--Shlosman-type condition $K(L,\beta,\gamma)$ at one scale $L$
gives uniqueness and decay at rate $-\ln\gamma/(2L)$. A finite-volume coupling
$z(L,\beta)=L\,E_f(L,\beta)$, with $E_f$ the lowest energy in one unit of 't~Hooft electric flux
in the periodic box, is claimed to grow with $L$ for $\SU(2)$. If $z$ became large at some scale
and large $z$ implied $K$, the gap would follow.

\begin{lemma}\label{neg:lem-decaygap}
Let $T=e^{-H}$ be the transfer matrix of a reflection-positive state, $0\le T\le1$. Suppose
that for a set of local observables $F$ whose vectors $\psi_F=(F-\langle F\rangle)\Omega$ are
dense in $\Omega^\perp$, $0\le\langle\psi_F,T^n\psi_F\rangle\le K_Fe^{-mn}$ for all $n\ge0$, with
$m$ independent of $F$. Then $\operatorname{spec}(H|_{\Omega^\perp})\subset[m,\infty)$.
\end{lemma}

\begin{proof}
Let $\nu_F$ be the spectral measure of $\psi_F$. For $m'<m$,
$e^{-nm'}\nu_F([0,m'))\le K_Fe^{-mn}$, so $\nu_F([0,m'))=0$. The spectral projection of
$[0,m')$ is bounded and vanishes on a dense set, hence on $\Omega^\perp$.
\end{proof}

The criterion itself (decay from $K$) is claimed: the Dobrushin--Shlosman paper was not read,
and its hypotheses for link variables were not checked.

\begin{observation}\label{neg:obs-fs}
The route fails at two missing steps.
\textup{(1)} For every $\beta>0$, $\sup_Lz(L,\beta)=\infty$. This follows from a positive string
tension at every $\beta$, which is open \cite{Ch21}; it is a confinement statement.
\textup{(2)} There are $z^*<\infty$, $k\in\N$ and $\gamma<1$ such that for every $\beta$ and $L$,
$z(L,\beta)\ge z^*$ implies $K(kL,\beta,\gamma)$. No inequality of any kind is known connecting
a periodic-box, one-sector number to a condition uniform over exterior configurations, even at
strong coupling.
The claimed monotonicity of $z$ rests on the sign of the one-loop coefficient
$b_0=11\cdot2/(48\pi^2)$ and on Table~\ref{neg:tab-fs}. Status: Monte Carlo by one program, one
$\beta$, three sizes, effective mass at $t=1$ (for $\SU(2)$, $m_{\mathrm{eff}}(2)<m_{\mathrm{eff}}(1)$,
so $z$ is overestimated). The sources do not give the seeds. The predicate ``$z_{\U(1)}$ within a
factor $2$ of $1/(2\beta)=0.217$ at $L=4,6$'' failed; that was the worker's expectation, which
ignored the finite temporal extent. A rotor model fitted afterwards reproduces the drift with
residuals of $1.5$ to $2.6$ standard errors.
\end{observation}

\begin{table}[ht]
\centering
\small
\begin{tabular}{@{}llll@{}}
\toprule
$L$ & $z$, $\SU(2)$ & $z$, $\U(1)$ & rotor model at $T=12$ (post hoc) \\
\midrule
$2$ & $0.562\pm0.016$ & $0.335\pm0.024$ & $0.324$ \\
$4$ & $2.12\pm0.06$ & $0.553\pm0.078$ & $0.422$ \\
$6$ & $8.02\pm0.40$ & $0.684\pm0.068$ & $0.500$ \\
\bottomrule
\end{tabular}
\caption{Flux coupling $z=L\,m_{\mathrm{eff}}(1)$ on $L^3\times12$ at $\beta=2.3$, from the
correlator of the zero-momentum spatial Polyakov loop. Free-photon value $1/(2\beta)=0.217$.}
\label{neg:tab-fs}
\end{table}

The table shows a steep rise for $\SU(2)$ and a slow drift for $\U(1)$ that is probably thermal.
It shows confinement-scale physics at $\beta=2.3$, not running from weak coupling, which would
need $L$ of order $\exp(\mathrm{const}\cdot\beta)$.

\emph{Discriminators.} For $\U(1)$ the route stops at the growth of $z$, as it must; for finite
groups $z\to0$ in the deconfined phase while (A) holds, so the route cannot reach them.

\subsection{Positivity of the character expansion}\label{neg:sub-dual}

\emph{Mechanism.} Expand each plaquette weight in characters and integrate out the links. If
every resulting amplitude were nonnegative, $\SU(2)$ would be a positive random-surface model.
Percolation-type sharpness arguments could then give the gap without any appeal to a running
coupling.

\begin{lemma}\label{neg:lem-cj}
For $\beta>0$ and $j\in\frac12\N_0$,
$e^{\frac\beta2\Tr U}=\sum_jc_j(\beta)\chi_j(U)$ with
$c_j(\beta)=I_{2j}(\beta)-I_{2j+2}(\beta)=\frac{2(2j+1)}{\beta}I_{2j+1}(\beta)>0$.
\end{lemma}

\begin{proof}
For $U$ with eigenvalues $e^{\pm i\theta}$, $\frac12\Tr U=\cos\theta$ and
$\chi_j=\sin((2j+1)\theta)/\sin\theta$. From $e^{\beta\cos\theta}=\sum_{n\in\Z}I_n(\beta)e^{in\theta}$,
the coefficient of $\sin((2j+1)\theta)$ in $e^{\beta\cos\theta}\sin\theta$ is $I_{2j}-I_{2j+2}$.
The recurrence $I_{\nu-1}-I_{\nu+1}=(2\nu/\beta)I_\nu$ gives the second form, and $I_n(\beta)>0$
because its power series has positive terms.
\end{proof}

The identity was also checked in floating point for $2j\le10$ and five values of $\beta$
(relative difference below $10^{-15}$). For a configuration $\{j_p\}$ of plaquette spins put
$A(\{j_p\})=\int\prod_p\chi_{j_p}(U_p)\prod_e dU_e$. Then
$Z=\sum_{\{j_p\}}\prod_pc_{j_p}(\beta)\,A(\{j_p\})$. For an abelian group $\prod_p\chi_{j_p}(U_p)$
is a character of the group of link variables, so $A\in\{0,1\}$.

\begin{observation}\label{neg:obs-dual}
Not every $\SU(2)$ amplitude is nonnegative. On the periodic $2\times1\times1$ lattice in three
dimensions, $4$ of the $729$ spin configurations considered have $A<0$, $141$ have $A>0$ and
$584$ have $A=0$. The minimum is $A=-0.003086419753$, at spins $(j_{01},j_{02},j_{12})=(0,1,1)$ on
site $0$ and $(1,0,\frac12)$ on site $1$. A separate Monte Carlo program gives
$-0.002892\pm0.000237$ (seed $7$, $4\cdot10^7$ samples), $12$ standard errors below zero.
Setting every plaquette that touches a fourth direction to spin $0$ embeds the configuration in
four dimensions with the same amplitude (proved by the worker: those links occur in no character).
Status: computed in floating point by one program (intertwiner contraction), confirmed by a
second program by a different method. The value agrees with $-1/324$ to the digits printed;
the rational value is not proved. On the one-site three-torus all $27$ amplitudes are positive,
so the smallest lattice does not show the sign.
\end{observation}

The literature knows this as the sign problem of the dual \cite{CCK07}; that source was read in
search snippets only. The test refutes positivity of the plaquette-character dual only. It does
not exclude a positive representation obtained by resummation over intertwiners, or any other
representation.

\emph{Discriminators.} For $\U(1)$ and $\Z_2$ the amplitudes are indicators and positivity holds;
since $\U(1)$ in four dimensions has no gap at weak coupling, positivity alone could never have
carried the route, and for $\Z_2$ the surface model has a transition at large $\beta$ while (A)
holds.

\subsection{Exclusion}\label{neg:sub-excl}

\emph{Mechanism.} $\SU(2)$ has no continuous global or one-form symmetry, only its center
$\Z_2$. If ``no continuous symmetry to break'' implied a gap except at isolated couplings, (A)
would reduce to excluding a transition at weak coupling. That exclusion would follow if the
zero set of $m$ were closed, the correlation length diverged at its infimum, and the free energy
were non-analytic there.

\begin{lemma}\label{neg:lem-usc}
Assume \textup{(H1)} a unique infinite-volume Gibbs state $\mu_\beta$ for every $\beta$;
\textup{(H2)} the specification is continuous in $\beta$ and in the boundary condition;
\textup{(H3)} reflection positivity, so that $C_F(n)=\langle\psi_F,T^n\psi_F\rangle$;
\textup{(H4)} vectors of local positive-time observables are dense. Then
$m(\beta)=\inf_F\inf_k\bigl(-\tfrac1{2k}\log(C_F(2k)/C_F(0))\bigr)$ is upper semicontinuous in
$\beta$. In particular $\{\beta:m(\beta)\ge c\}$ is closed for every $c$.
\end{lemma}

\begin{proof}
(H1), (H2) and compactness give weak continuity of $\beta\mapsto\mu_\beta$, so each $C_F(n)$ is
continuous. By the Lyapunov inequality, $-\frac1{2k}\log(C_F(2k)/C_F(0))$ is nonincreasing in $k$
with limit $-\log\|T|_{\overline{\mathrm{span}}}\|$. By (H4) the infimum over $F$ is $m$. An
infimum of continuous functions is upper semicontinuous.
\end{proof}

(H1) is the standing assumption of the paper and is not known at large $\beta$; (H2) is claimed
as standard, with no source read.

\begin{observation}\label{neg:obs-usc}
Under \textup{(H1)--(H4)} nothing beyond upper semicontinuity follows from continuity of the
correlations. An abstract continuous family of spectral measures has zero set
$\{m=0\}=(1,\infty)$, which is not closed, and a correlation length equal to $1$ up to the
transition and infinite after it (checked by one program). It is not a lattice model. No theorem
of the form ``the gap vanishes $\Rightarrow$ the free energy is non-analytic'' was found.
\end{observation}

The principle itself is false; the counterexamples are Observation~\ref{neg:obs-excl}.

\emph{Discriminators.} For $\U(1)$ the route is consistent but uninformative (it has a continuous
one-form symmetry and no gap); the finite groups $\Z_N$ with large $N$ are the counterexample to
the principle.

\section{A shared-properties observation}\label{sec:shared}

\begin{observation}\label{neg:obs-shared}
The following six properties hold for $\SU(2)$, $\U(1)$ and $\Z_2$ Wilson lattice gauge theory
in four dimensions.
\begin{enumerate}
\item[P1] Reflection positivity at every $\beta$. Cited \cite{OS78}, as stated in its abstract.
\item[P2] A strictly positive self-adjoint transfer matrix with a Perron--Frobenius ground
state. Positivity of the transfer matrix is cited \cite{OS78,Lu77}, as stated in their
abstracts; the Perron--Frobenius ground state is argued in the worker report from the
Hamiltonian limit, not cited.
\item[P3] Positive character coefficients of the Wilson weight at every $\beta>0$. Proved here
for $\SU(2)$ (Lemma~\ref{neg:lem-cj}); for $\U(1)$ the coefficients are $I_n(\beta)>0$; for
$\Z_2$ they are $\cosh\beta$ and $\sinh\beta$.
\item[P4] A gap and analyticity at small $\beta$. Cited \cite{OS78} for the gap and area law at
strong coupling for every compact group, as stated in its abstract; analyticity by cluster
expansion is claimed as standard.
\item[P5] A positive gap in every finite periodic volume at every $\beta$. Stated in the worker
reports, not written out there; it follows from P2 by the Perron--Frobenius theorem for the
compact finite-volume transfer operator.
\item[P6] A nontrivial center: $\Z_2\subset\SU(2)$; $\U(1)$ and $\Z_2$ are their own centers.
Elementary.
\end{enumerate}
Compact $\U(1)$ in four dimensions has P1--P6 and no gap at weak coupling \cite{Guth80,FS82}
(as stated in their abstracts). Consequently no argument that uses only P1--P6 proves (A).
\end{observation}

The observation does not exclude an argument that uses a property outside the list. It does
not exclude an argument that uses the dimension in a way that fails for $\U(1)$ in four
dimensions, and it does not exclude an argument that uses a property $\U(1)$ lacks. Three such
properties are known with proof or computation:
\begin{enumerate}
\item $\SU(2)$ is perfect: $[\SU(2),\SU(2)]=\SU(2)$, so it has no nontrivial homomorphism to
$\U(1)$. Elementary.
\item Its Ricci curvature is the positive constant $1$, against $0$ for $\U(1)$
(Proposition~\ref{neg:prop-flat}). Proved.
\item Its plaquette-character amplitudes can be negative, while those of $\U(1)$ and $\Z_2$ are
indicators (Observation~\ref{neg:obs-dual}). Computed in floating point and confirmed by Monte
Carlo; not proved in exact arithmetic.
\end{enumerate}
All three are local. None has a known global consequence for the gap. The official problem
statement records the same absence: there is no known idea for a gap uniform in the volume
\cite{JW}. In the routes of
Section~\ref{sec:routes} the second enters only at flat configurations and is outweighed on
typical ones (Observation~\ref{neg:obs-curv}); the third enters as cancellation, which removes
the positivity tools. Three other candidates meet the same wall. Analyticity in $\beta$ is
proved only at small $\beta$, and at all $\beta$ it would exclude a transition that $\U(1)$ has.
A uniform gap of the Langevin generator is the strong mixing of
Observation~\ref{neg:obs-curv-missing}. No inequality monotone in the dimension is known, and
$\U(1)$ has a gap in three dimensions \cite{GM82} (as stated in its abstract) and none in four.

\begin{observation}\label{neg:obs-excl}
The principle ``a lattice gauge theory with no continuous global or one-form symmetry has a
gap except at isolated couplings'' fails.
\begin{enumerate}
\item Four-dimensional $\Z_N$ lattice gauge theory with $N$ large has an intermediate massless
phase, on an interval of couplings, in which local $\U(1)$ invariance is restored dynamically
\cite{FS82}. Status: cited, as stated in its abstract. The action, the exact form of the
massless statement and the interval were not read. $\Z_N$ has only a $\Z_N$ one-form symmetry.
\item Five-dimensional $\SU(2)$ Wilson theory has a first-order bulk transition at
$\beta_c\approx1.64$ into a phase with zero string tension \cite[\S3 and Fig.~1]{FiveD11},
credited there to \cite{Cr79} (as stated in its abstract). Status: numerical. That the
weak-coupling phase is massless is perturbative expectation; it was not found stated in the
text read, and it is not proved. A Monte Carlo by one program on $4^5$ (200 Metropolis sweeps,
one seed per start, written by a sub-worker and rerun bit-identical by the worker) gives a
hot--cold hysteresis gap of $0.090\pm0.003$ at $\beta=1.6$, $0.018$ at $\beta=1.7$ (still
drifting), and below $0.0015$ at $\beta=1.3$, $1.5$, $1.9$ and $2.2$ (the same report's summary
table says below $0.002$). Its weak-coupling plaquette approaches the one-loop value
$1-3/(5\beta)$: $0.6925$ against $0.7273$ at $\beta=2.2$ and $0.630$ against $0.684$ at
$\beta=1.9$. This is consistent with a perturbative phase and is not a measurement of a massless
correlator.
\end{enumerate}
$\SU(2)$ in five dimensions is compact, connected, simple and perfect, has positive Ricci
curvature and signed amplitudes, and its reflection positivity in site planes does not depend on
the dimension (claimed, re-derived for four dimensions only). So an exclusion argument for (A)
must use $d=4$ in a way that fails at $d=5$, and must forbid an emergent continuous one-form
symmetry, as in the $\Z_N$ example. Both are dynamical statements about the running of the
coupling.
\end{observation}

\section{Hypotheses (C) and (E)}\label{sec:CE}

\subsection{Hypothesis (C)}

Let $T=e^{-H}$ be the transfer matrix at $\beta$, $\psi\perp\Omega$ the vector of a local
observable, $\nu$ its spectral measure, $C(n)=\langle\psi,T^n\psi\rangle=\int e^{-n\varepsilon}\,d\nu(\varepsilon)$
for real $n\ge0$, $m=m(\beta)$ and $\xi=1/m$. We write (C) in the form
$C(\xi/2)\le D\,C(\xi)$ with $D$ independent of $\beta\ge\beta_1$. The effective mass
$u(n)=-\frac{d}{dn}\log C(n)$ is the mean energy of $e^{-nH/2}\psi$; (C) is equivalent, up to
constants and a change of time from $\xi/2$ to $\xi$, to $u(\xi/2)\le K\,m$.

\begin{proposition}\label{neg:prop-Cprime}
Assume \textup{(C$'$)}: $1-C_\beta(1)/C_\beta(0)\le K'm(\beta)$ for all $\beta\ge\beta_1$. Then for
every $j\ge0$ with $\xi/2^j\ge1$ and $\xi\ge2K'$,
$C(\xi/2^j)\le e^{2K'}C(\xi)$. In particular (C) holds with $D=e^{2K'}$.
\end{proposition}

\begin{proof}
$(\log C)''$ is the variance of $\varepsilon$ under $e^{-n\varepsilon}d\nu/C(n)$, so $\log C$ is
convex. Put $q=C(1)/C(0)$. For $a\ge1$ the slope of $-\log C$ on $[a,\infty)$ is at most its
chord slope on $[0,1]$, which is $\log(1/q)$. Hence
$\log C(\xi/2^j)-\log C(\xi)\le\xi\log(1/q)\le2\xi K'm=2K'$, using $-\log(1-x)\le2x$ for
$x\le\frac12$.
\end{proof}

The worker compiled the integer form of this chain (Lean, declaration \texttt{A3C.step\_chain})
with exit code $0$ and the axioms propext, Classical.choice, Quot.sound; it was not recompiled
for this section. The real-time convexity argument is not formalized.

\begin{observation}\label{neg:obs-Cprime-free}
(C$'$) fails for lattice-local observables in the free massive lattice field, while (C) holds
there. For $\varphi(0)^2$, $\xi(1-C(1)/C(0))=0.963\,\xi$ for $\xi$ from $4$ to $256$, and $C(1)/C(0)$
tends to the lattice constant $0.0373$. For $\varphi(0)$, $R(\xi)=C(\xi/2)/C(\xi)$ is $6.584$,
$5.921$, $5.580$, $5.522$, $5.508$, $5.505$, $5.504$ at $\xi=4,8,\dots,256$, against the continuum
value $2K_1(1/2)/K_1(1)=5.503975$; for $\varphi(0)^2$ it tends to $30.2937$. Status: computed in
floating point by one program (Schwinger representation); the values of $R$ for $\varphi(0)$ were
recomputed by a second program (momentum integral), agreeing to $10^{-6}$ relative. So (C$'$)
can be asked only of an observable smeared on the physical scale. For the slice-flowed free
field at $s=0.1$, $\xi(1-C(1)/C(0))=3.56$, $4.13$, $4.46$, $4.64$, $4.73$, $4.78$ at
$\xi=8,16,32,64,128,256$ (one program): bounded, as (C$'$) requires.
\end{observation}

What (C) needs for $\SU(2)$ is therefore missing: (C$'$) for the slice-flowed magnetic density at
flow time $s\xi(\beta)^2$, in infinite volume, uniformly in $\beta$.

\begin{proposition}\label{neg:prop-smear}
Smearing in Euclidean time by a fixed number of correlation lengths does not make (C)
automatic. Let $k\ge1$, $\psi=T^{k\xi/2}\varphi$, and let $\varphi$ have spectral measure
$\nu_\varphi=w\,\delta_m+\delta_1$ with $m=1/\xi<1$ and $w=e^{-(k+1)(\xi-1)}$. Then
$C_\psi(\xi/2)/C_\psi(\xi)\ge\frac12e^{\xi/2}$.
\end{proposition}

\begin{proof}
$C_\psi(n)=C_\varphi(n+k\xi)$ and $C_\varphi(n)=we^{-n/\xi}+e^{-n}$. So
$C_\varphi((k+1)\xi)=2e^{-(k+1)\xi}$ and $C_\varphi((k+\frac12)\xi)\ge e^{-(k+\frac12)\xi}$.
\end{proof}

Smearing by a $\beta$-dependent time $n_\beta$ chosen so that $u_\varphi(n_\beta)\le(1+\delta)m$
does give (C), trivially. If $n_\beta/\xi$ is unbounded, the vector is no longer a sharp-time
local field; if it is bounded, that boundedness is already (C)-type information. By explicit
examples, (C) fails when the spectral weight at the bottom is exponentially small in $\xi$, a
hole in the spectral measure alone does not break it, and (C) does not exclude a limit that is a
single atom.

\begin{observation}\label{neg:obs-matched}
In a three-spacing $\SU(2)$ measurement of the slice-flowed magnetic density
(Table~\ref{neg:tab-matched}), $\xi(1-C(1)/C(0))$ is $0.53$, $0.68$, $0.86$ at $\xi=0.72$, $1.01$,
$1.40$. It increases, roughly like $\xi^{0.75}$, over the factor $1.9$ available. Status: Monte
Carlo by one set of programs, measured only. Caveats: the correlation lengths are about one lattice
step, the worst place to ask the question; $\xi$ is taken from a published glueball mass, not
measured; each of the two added seeds at $N=12$ repeats three configurations (IDs $40$--$42$)
after a restart, found by another director's audit, and the table gives the original
row-based estimates (according to the ledger, with the repeats removed the estimates move
within their errors); the boxes are matched in physical size only to $3$ to $7$ per cent.
\end{observation}

\begin{table}[ht]
\centering
\small
\begin{tabular}{@{}lllllllll@{}}
\toprule
$N^4$ & $\beta$ & $s/a^2$ & configs & $a\sqrt\sigma$ & $R(0)\cdot10^2$ & $R(1)\cdot10^2$ & $\xi$ & $\xi(1-q)$ \\
\midrule
$6^4$ & $2.3$ & $0.5$ & $2990$ & $0.3690(30)$ & $3.751(42)$ & $1.012(32)$ & $0.72$ & $0.53$ \\
$8^4$ & $2.4$ & $0.889$ & $1500$ & $0.2660(20)$ & $3.615(60)$ & $1.166(51)$ & $1.01$ & $0.68$ \\
$12^4$ & $2.5$ & $2.0$ & $840$ ($834$) & $0.1905(8)$ & $3.233(61)$ & $1.245(51)$ & $1.40$ & $0.86$ \\
\bottomrule
\end{tabular}
\caption{Slice-flowed magnetic density, flow to $s=\ell^2/72$ with $\ell$ the box side.
$R(\tau)$ is the connected correlator of the slice average over $\langle\bar E\rangle^2$, so
$q=R(1)/R(0)=C(1)/C(0)$: $0.270$, $0.323$, $0.385$. Scale $a\sqrt\sigma$ from \cite{FHK93}
(Table~3); $\xi=1/(3.74\,a\sqrt\sigma)$ with the scalar glueball mass $3.74(12)\sqrt\sigma$ of
\cite{Te98} (\cite{LTW04} gives $3.78(7)$). At $N=12$: three seeds, $840$ rows, $834$ distinct
configurations. The last three columns are derived; errors are not propagated because the
correlation of $R(0)$ and $R(1)$ is not given. The sources do not give the production seeds.}
\label{neg:tab-matched}
\end{table}

The measurement supports one sentence: on these lattices $\xi(1-C(1)/C(0))$ grows with $\xi$,
which is the direction that (C$'$) for the slice-flowed density does not predict. It refutes
neither (C$'$) nor (C).

\subsection{Hypothesis (E)}

(E) asks that the scaling limit be interacting and asymptotically free at short distance. A
generalized free field (GFF) is a Gaussian field with an arbitrary K\"all\'en--Lehmann measure;
a free algebra is the Wick-polynomial algebra of finitely many GFFs.

\begin{proposition}\label{neg:prop-GFF}
\textup{(i)} Every Euclidean-invariant, reflection-positive two-point function is the two-point
function of a GFF. Hence no statement about a single two-point function excludes a GFF.
\textup{(ii)} Let $O=\frac12\sum_{a=1}^k\varphi_a^2$ with $\varphi_a$ independent copies of one
GFF. Then the normalized fourth cumulant of $O$ at one point is $\kappa_4/\kappa_2^2=12/k$, which
is $4$ for $k=3$. Hence a nonvanishing truncated four-point function of a composite field does
not exclude a free algebra.
\end{proposition}

\begin{proof}
(i) Take the Gaussian measure with that covariance; a Gaussian measure is reflection positive
when its covariance is. (ii) $2O$ is $\sigma^2\chi^2_k$, whose cumulants are
$\kappa_n=2^{n-1}(n-1)!\,k$ in units of $\sigma^{2n}$; so $\kappa_2=2k$, $\kappa_4=48k$.
\end{proof}

Part (ii) applies to the rescaled plaquette energy used in the program's proof of (B),
Theorem~\ref{thm:main}. At fixed lattice distance and
$\beta\to\infty$ the limit of $\beta(1-\cos\alpha_p)$ is $\frac12\sum_{a=1}^3(F^a_p)^2$ (a claimed
step); if the three colour components are independent with equal variance (not checked), it
has $\kappa_4/\kappa_2^2=4$. A test ``$\kappa_4\ne0\Rightarrow$ interacting'' applied to it would
declare the free Maxwell limit interacting. A second mass, or any mass ratio, does not
distinguish either: a GFF with measure $\delta_{m_1^2}+\delta_{m_2^2}$ has any ratio.

A formulation that does exclude every free algebra is the two-particle energy shift. Assume
that the lightest mass $m$ is an isolated eigenvalue of the mass operator. Let $E_2(L)$ be the
lowest two-particle energy at zero total momentum in a periodic box of side $L=z/m$, and
$\delta(z)=\lim_{\beta\to\infty}(E_2(L)-2m)/m$. Then (E$^*$): $\delta(z)\ne0$ for some $z$ in some
symmetry channel. In the Fock space of a free algebra the Hamiltonian is a sum of one-particle
Hamiltonians, so $E_2(L)=2m$ exactly; hence (E$^*$) excludes every free algebra. It needs the
isolated one-particle shell, which a gap does not give. Its relation to the scattering length
through \cite{LW86} is claimed from memory; the source was not read.

\begin{table}[ht]
\centering
\small
\begin{tabular}{@{}lll@{}}
\toprule
model & quantity & values \\
\midrule
free field, $d=4$, $L=8$, $m=0.25$ & $U_4$; $g_R$; block $\kappa_4/\kappa_2^2$ & $0.007\pm0.025$; $0.34\pm1.20$; $-0.19\pm0.07$ \\
Ising, $d=4$, $\xi_2/L=0.25$ & $g_R$ at $L=6,8,12$ & $67\pm9$, $46\pm8$, $32\pm7$ \\
Ising, $d=4$, $\xi_2/L\approx0.33$ & $g_R$ at $L=6,8,12,16$ & $36\pm2$, $\approx29$, $22\pm2$, $20\pm3$ \\
Ising, $d=2$, $K_c$ & $U_4$ at $L=16,32$ & $0.6111\pm0.0005$, $0.6098\pm0.0007$ \\
\bottomrule
\end{tabular}
\caption{Calibration on known answers. Free field: $4000$ exact Gaussian samples, $\xi_2=4.023\pm0.046$
against the exact $4$. Four-dimensional Ising: Wolff cluster algorithm, $17$ runs, seeds
$2001$--$2017$, $15000$ measurements each. Two-dimensional Ising: seeds $31$, $32$; a separate
program (the worker's child) gives $0.6115$, $0.6107$. Monte Carlo, measured only. The
renormalized coupling is $g_R=3U_4(L/\xi_2)^4$ with $U_4=1-\langle M^4\rangle/(3\langle M^2\rangle^2)$.}
\label{neg:tab-ising}
\end{table}

\begin{observation}\label{neg:obs-ising}
On the calibration of Table~\ref{neg:tab-ising}, the free field gives $U_4$ and $g_R$ compatible
with zero (the block cumulant sits $3$ standard errors from zero and was not rerun with a second
seed). In four-dimensional Ising, $g_R$ at fixed $\xi_2/L$ decreases with $L$, consistent with
the Gaussian scaling limit of \cite{ADC21} (as stated in its abstract); a logarithmic law was not
tested, and the predicate ``$g_R(12)$ exceeds zero by five standard errors at $\xi_2/L=0.25$''
failed at $4.6$. In two-dimensional Ising at criticality $U_4\approx0.61$, far from zero. The
critical $U_4$ on a periodic torus is nonzero also for models with Gaussian limits, because of
the zero mode, so only $g_R$ at fixed $L/\xi_2$ calibrates the test.
\end{observation}

What is not known. No rigorous non-Gaussian scaling limit of a lattice gauge theory in three or
more dimensions is known to the workers; no search establishing absence was made. The constant
in the short-distance statement of (E) for the trace of the stress tensor is not derived. Whether
(E$^*$) can be proved from lattice information is open: it lives at the correlation length,
where control is the gap problem itself.
